\section{引言：为什么关注 MoE？}

本期是「大语言模型架构」系列的第一讲。近年来，大模型架构的演进日新月异，其中 \textbf{Mixture of Experts (MoE)} 架构因其在推理效率上的巨大优势而备受瞩目。

以 DeepSeek R1 为例：它是一个约 670B 总参数的 MoE 模型，但推理时仅激活约 72B 的参数。与一个完整的 72B Dense 模型相比，两者的推理速度基本一致——这意味着 MoE 模型在保持巨大模型容量（能力强）的同时，推理效率极高（激活参数少）。

\begin{importantbox}{MoE 的核心价值}
\begin{itemize}
  \item \textbf{总参数量大} $\Rightarrow$ 模型容量大、能力强
  \item \textbf{激活参数量小} $\Rightarrow$ 推理速度快、效率高
  \item Qwen2.5 时代还是全 Dense 模型，到 Qwen3 已大量采用 MoE 变体
\end{itemize}
\end{importantbox}

本期以 \textbf{Qwen3-32B}（Dense）和 \textbf{Qwen3-30B-A3B}（MoE）为对象，通过参数量的对比计算来理解 MoE 架构。

